% ch1_a §1.1–1.2 (书页 1–8 / p015–p022)
% 第1章 PERIODIC STRUCTURES

\chapter{周期结构}

\epigraphbook{再来一遍！再来一遍！再来一遍！}{Thomas Campbell}

\section{平移对称性}

如果大多数固体的最稳定结构不是规则的晶格，那么关于固体物理性质的理论实际上将无从建立。平移对称性（translational symmetry）的存在，使 $N$ 体问题被约化到可以处理的规模。这就是说，\emph{存在着基矢（basic vectors）$\mathbf{a}_1$、$\mathbf{a}_2$、$\mathbf{a}_3$，使得原子结构在经过任何一个矢量——它为这些基矢的整数倍之和——的平移后保持不变}。

\emph{实际上}，这只是一个理想。实验室样品必然尺寸有限，因此我们不能把结构一直延伸越过边界。但受此影响的区域只有表面附近的几层原子；在一块含 $N$ 个原子的样品中，它们只占大约 $N^{2/3}$ 个原子——就是说，在宏观样品中不过 $10^{8}$ 个原子里有一个。大多数晶体固体在结构上也是不完美的，存在缺陷、杂质和位错，扰乱着原子排列的规则性。这类不完整性引起许多有趣的物理现象，但在本书的讨论中，除非偶然涉及，我们将一概置之不理。我们在这里主要关注的，是完美的理想固体及其所表现的性质；而那些把固体仅仅当作些许污垢、微细裂纹和结构瑕疵的基体、载体或背景的现象，则属于另一个论域。

我们用空间点阵（space lattice）或 Bravais 网（Bravais net）来表示平移群。从某一点出发，构造出由基本平移可从它到达的所有点。这些点便是格点（lattice sites），由集合
\begin{equation*}
\boldsymbol{l} = l_1\mathbf{a}_1 + l_2\mathbf{a}_2 + l_3\mathbf{a}_3,
\tag{1.1}
\end{equation*}
给出，其中 $l_1$、$l_2$、$l_3$ 为整数（图 1）。

但固体是一种物理结构——并不是一组数学点。设想在我们所选定的原点 $O$ 附近有若干原子等等。平移不变性坚持要求：在每个格点周围，都必须恰好有完全相同的原子，以完全相同的方式配置（图 2）。

显然，只要指明单个晶胞（unit cell）的内容——例如由基矢 $\mathbf{a}_1$、$\mathbf{a}_2$、$\mathbf{a}_3$ 张成的平行六面体——我们就能确定整个晶体的物理布局。整个晶体就是这种物体的无穷多次重复，像砌墙的砖一样堆叠起来。不过，晶胞的具体取法在某种程度上是任意的：任何大小、形状与取向合适的平行六面体显然都行——正如我们在
%FIG{1}: {图 1}
%FIG{2}: {图 2. 可供选择的晶胞。}
图 2 中看到的那样。或许不那么显而易见的是，晶胞的形状还可以在一定程度上加以改变。例如，假设结构中某一点具有中心对称性（因而在所有等价点处也都是如此）。这一点便很适合选作晶胞的中心，晶胞自身也具有中心对称性。系统地做到这一点的办法，是构造一个 Wigner–Seitz 原胞（Wigner–Seitz cell），也就是从选定的中心向最近邻的等价格点，画出各平移矢量的垂直平分面。全部平分面之内的体积显然是一个晶胞——它就是其中各点离选定中心比离任何其他格点都近的那个区域。

晶胞可以含有一个或多个原子。自然，如果只含一个原子，我们就把它置于格点中心，并说我们有一个 Bravais 格子（Bravais lattice）。如果每个晶胞含几个原子，那么我们就有了一个带基元的格子（lattice with a basis）。在后面的大部分内容中，我们将不加特别声明地假定结构是 Bravais 格子。这只是为简单起见；实际上只有少数几种单质固体——例如碱金属——才具有这种结构。

晶体学（crystallography）这门科学，关注的是对所有可能类型的晶体结构进行清点和分类，并确定实际晶体的实际结构。

%FIG{3}: {图 3. Wigner–Seitz 原胞。}

结构按其对称性质分类，例如绕轴旋转的不变性、对平面反射的不变性等等。这些对称性常常对简化理论计算十分重要，并且在讨论定义固体宏观性质所需要的参数数目时，可以发挥很大的威力。然而，要充分利用这套理论，需要群论（group theory）这门数学，它会使我们偏离正题太远。如果我们主要局限于非常简单的固体，那么大多数对称性质无须正式的代数分析，单凭观察就能发现。无论如何，关于晶体学、以及关于群论及其在固体理论中的应用，都有许多优秀的著作。

在这些书中，各类 Bravais 格子等等都有详细的铺陈。我们在这里只考虑一个例子，它体现了本学科的许多原理，而且作为一种某些元素实际上采取的结构，也极为重要。这就是体心立方（body-centred cubic，B.C.C.）结构，如图 4 所示。

%FIG{4}: {图 4. 体心立方格子。（a）立方晶胞。（b）Bravais 格子的生成矢。（c）Wigner–Seitz 原胞。}

%FIG{5}: {图 5. B.C.C. 格子 Wigner–Seitz 原胞的堆垛。}

乍一看，这是一个每个立方晶胞含两个原子的格子，或者说是两个相互穿插的简立方亚格子，由
\begin{equation*}
\left.
\begin{aligned}
\boldsymbol{l} &= l_x\mathbf{a}_x + l_y\mathbf{a}_y + l_z\mathbf{a}_z, \\
\boldsymbol{l}' &= \left(l_x + \tfrac{1}{2}\right)\mathbf{a}_x + \left(l_y + \tfrac{1}{2}\right)\mathbf{a}_y + \left(l_z + \tfrac{1}{2}\right)\mathbf{a}_z,
\end{aligned}
\right\}
\tag{1.2}
\end{equation*}
定义，其中 $l_x$、$l_y$、$l_z$ 均为整数。但如果写下
\begin{equation*}
\left.
\begin{aligned}
\mathbf{a}_1 &= \tfrac{1}{2}\left(-\mathbf{a}_x + \mathbf{a}_y + \mathbf{a}_z\right), \\
\mathbf{a}_2 &= \tfrac{1}{2}\left(\mathbf{a}_x - \mathbf{a}_y + \mathbf{a}_z\right), \\
\mathbf{a}_3 &= \tfrac{1}{2}\left(\mathbf{a}_x + \mathbf{a}_y - \mathbf{a}_z\right),
\end{aligned}
\right\}
\tag{1.3}
\end{equation*}
我们便能由
\begin{equation*}
\boldsymbol{l} = l_1\mathbf{a}_1 + l_2\mathbf{a}_2 + l_3\mathbf{a}_3
\tag{1.4}
\end{equation*}
（$l_1$、$l_2$、$l_3$ 为整数）生成两个亚格子的全部格点。我们发现，落在立方体的中心还是顶角上，视 $(l_1+l_2+l_3)$ 为奇数还是偶数而定。

因此（图 4(b)），这其实是一个 Bravais 格子。我们可以不用立方晶胞，而改用 Wigner–Seitz 原胞（图 4(c)）；它的构造方法是，沿从中心到角点的对角线在半程处把立方体的各个角截去。这个图形显然具有与立方体相同的对称性——例如，原来的矢量 $\mathbf{a}_x$、$\mathbf{a}_y$、$\mathbf{a}_z$ 都是四重对称轴。它还（或许比原来的立方体更清楚地）表明，矢量 $\mathbf{a}_1$、$\mathbf{a}_2$、$\mathbf{a}_3$，即立方体的对角线，是三重旋转对称轴，因为它们穿过原胞的六边形面。设想一下这些多面体如何能堆砌成原来的格子（图 5），是一次很好的练习。

\section{周期函数}

要为一个晶体结构定义物理模型，我们需要给空间中的某个函数 $f(\mathbf{r})$ 赋值——局域电子密度、静电势等等，它们应能被辨认为原子的一种排列（例如图 6）。我们关于平移对称性的断言意味着，这必须是一个\emph{多重周期函数}（multiply periodic function）
\begin{equation*}
f(\mathbf{r}+\boldsymbol{l}) = f(\mathbf{r})
\tag{1.5}
\end{equation*}
对空间中所有的点 $\mathbf{r}$、对所有的格平移（lattice translations）都成立。

%FIG{6}: {图 6. 多重周期函数。}

%FIG{7}: {图 7}

在一维情形，我们对周期函数再熟悉不过了。例如，如图 7 那样，可以有
\begin{equation*}
f(x+l) = f(x),
\tag{1.6}
\end{equation*}
其中 $l$ 取 $l_1 a$ 的形式，$l_1$ 为整数，而 $a$ 是函数的周期。我们还知道，任何这类函数都可以表示成 Fourier 级数，
\begin{equation*}
f(x) = \sum_n A_n e^{2\pi i n x/a},
\tag{1.7}
\end{equation*}
其中 $n$ 是整数。让我们把它写成
\begin{equation*}
f(x) = \sum_g A_g e^{igx},
\tag{1.8}
\end{equation*}
其中的量 $g$ 属于倒格子长度（reciprocal lattice lengths）\footnote{在这个定义中把因子 $2\pi$ 包括进来是方便的；虽然在晶体学中惯例是把 $2\pi$ 明确地留在指数上。}的集合
\begin{equation*}
g_n = n\,\frac{2\pi}{a}.
\tag{1.9}
\end{equation*}

大家熟知，(1.8) 中的系数由
\begin{equation*}
A_g = \frac{1}{a}\int_{\text{cell}} f(x)\, e^{-igx}\, dx,
\tag{1.10}
\end{equation*}
定义；在这个例子里，积分范围譬如取 $0 < x \leqslant a$——它只需是晶格的一个晶胞。(1.8) 蕴含 (1.6) 这一初等证明，可以由如下条件导出：对任何 $g$ 值，
\begin{equation*}
e^{igl} = 1
\tag{1.11}
\end{equation*}
对所有平移 $l$ 成立。当 $g$ 取其一个允许值 $g_n$ 而 $l = l_1 a$ 时，这不过是说
\begin{equation*}
gl = n\,\frac{2\pi}{a}\, l_1 a = n l_1\, 2\pi = \text{integer} \times 2\pi
\tag{1.12}
\end{equation*}
罢了。(1.10) 的推导同样可以由这同一个关系得出。我们在这里不关心数学上病态的函数，尽可以放心地随意使用朴素的 Fourier 理论。

把这个定理推广到具有三条正交轴的结构相当容易。设这些轴为 $\mathbf{a}_x$、$\mathbf{a}_y$、$\mathbf{a}_z$，那么，只要
\begin{equation*}
f(\mathbf{r}+\boldsymbol{l}) \equiv f(\mathbf{r}+l_x\mathbf{a}_x+l_y\mathbf{a}_y+l_z\mathbf{a}_z) = f(\mathbf{r}),
\tag{1.13}
\end{equation*}
便有
\begin{equation*}
f(\mathbf{r}) = \sum_{g_x,g_y,g_z} A(g_x,g_y,g_z) \exp\left\{ i(g_x x + g_y y + g_z z)\right\},
\tag{1.14}
\end{equation*}
其中 $g_x$ 等等各自都是从集合 $2\pi n/a_x$ 等等中取的一个倒格子长度。证明可以这样做：先把 $f(\mathbf{r})$ 沿 $x$ 方向展开成 Fourier 级数，然后证明该级数的每个系数都是 $y$ 的周期函数，因而可以再展开成一个级数——对 $z$ 亦照此办理。

让我们把 (1.14) 改写成
\begin{equation*}
f(\mathbf{r}) = \sum_{\boldsymbol{g}} A_{\boldsymbol{g}} e^{i\boldsymbol{g}\cdot\mathbf{r}},
\tag{1.15}
\end{equation*}
其中 $\boldsymbol{g}$ 是分量为 $(g_x,g_y,g_z)$ 的矢量。这种矢量具有 (1.11) 所表达的性质，也就是说，对任何一个这样的矢量
\begin{align*}
\boldsymbol{g}\cdot\boldsymbol{l} &= \left(g_x l_x \mathbf{a}_x + g_y l_y \mathbf{a}_y + g_z l_z \mathbf{a}_z\right) \\
&= \frac{2\pi n_x}{a_x} l_x \mathbf{a}_x + \frac{2\pi n_y}{a_y} l_y \mathbf{a}_y + \frac{2\pi n_z}{a_z} l_z \mathbf{a}_z \\
&= 2\pi \times \text{integer},
\tag{1.16}
\end{align*}
无论 $\boldsymbol{l}$ 取什么值。于是，
\begin{equation*}
e^{i\boldsymbol{g}\cdot\boldsymbol{l}} = 1
\tag{1.17}
\end{equation*}
对一切格矢 $\boldsymbol{l}$、对一切倒格矢（reciprocal lattice vector）$\boldsymbol{g}$ 都成立。

显然，这个条件足以使级数 (1.15) 表示一个像 (1.5) 那样具有晶格周期性的函数：
\begin{align*}
f(\mathbf{r}+\boldsymbol{l}) &= \sum_{\boldsymbol{g}} A_{\boldsymbol{g}}\, e^{i\boldsymbol{g}\cdot(\mathbf{r}+\boldsymbol{l})} = \sum_{\boldsymbol{g}} A_{\boldsymbol{g}}\, e^{i\boldsymbol{g}\cdot\mathbf{r}}\, e^{i\boldsymbol{g}\cdot\boldsymbol{l}} \\
&= \sum_{\boldsymbol{g}} A_{\boldsymbol{g}}\, e^{i\boldsymbol{g}\cdot\mathbf{r}} = f(\mathbf{r}).
\tag{1.18}
\end{align*}
人们不难设计出一个证明，说明它也是必要条件，也就是说，级数 (1.15) 中只能含有 $g$ 值满足条件 (1.17) 的那些项。

剩下的只是为非正交格子构造倒格矢。这很容易做到：取基矢的倒格子三矢组（reciprocal triad），即
\begin{equation*}
\mathbf{b}_1 = \frac{\mathbf{a}_2\wedge\mathbf{a}_3}{\mathbf{a}_1\cdot\mathbf{a}_2\wedge\mathbf{a}_3}, \qquad
\mathbf{b}_2 = \frac{\mathbf{a}_3\wedge\mathbf{a}_1}{\mathbf{a}_1\cdot\mathbf{a}_2\wedge\mathbf{a}_3}, \qquad
\mathbf{b}_3 = \frac{\mathbf{a}_1\wedge\mathbf{a}_2}{\mathbf{a}_1\cdot\mathbf{a}_2\wedge\mathbf{a}_3},
\tag{1.19}
\end{equation*}
并写出
\begin{align*}
\boldsymbol{g} &= g_1\mathbf{b}_1 + g_2\mathbf{b}_2 + g_3\mathbf{b}_3 \\
&= 2\pi n_1\mathbf{b}_1 + 2\pi n_2\mathbf{b}_2 + 2\pi n_3\mathbf{b}_3,
\tag{1.20}
\end{align*}
其中 $n_1$ 等等为整数。

由初等矢量分析即可验证：$\mathbf{b}_1\cdot\mathbf{a}_1 = 1$ 等等，以及 $\mathbf{b}_1\cdot\mathbf{a}_2 = 0$ 等等；于是
\begin{align*}
\boldsymbol{g}\cdot\boldsymbol{l} &= \left(g_1\mathbf{b}_1 + g_2\mathbf{b}_2 + g_3\mathbf{b}_3\right)\cdot\left(l_1\mathbf{a}_1 + l_2\mathbf{a}_2 + l_3\mathbf{a}_3\right) \\
&= g_1 l_1 + g_2 l_2 + g_3 l_3 \\
&= 2\pi \times \text{integer}.
\tag{1.21}
\end{align*}
因此，集合 (1.20) 中的每一个矢量都满足条件 (1.17)。

我们还可以把公式 (1.10) 加以推广，用来求级数 (1.15) 中的每个系数。它必须是一个体积分，我们把它写作
\begin{equation*}
A_{\boldsymbol{g}} = \frac{1}{v_{\text{cell}}} \int_{\text{cell}} f(\mathbf{r})\, e^{-i\boldsymbol{g}\cdot\mathbf{r}}\,\mathrm{d}\mathbf{r}.
\tag{1.22}
\end{equation*}
